\section{Вывод}

В ходе работы экспериментально определён КНД пирамидальной рупорной антенны зеркальным методом (методом Парселла). Предварительно по зависимости $|E|^2(x)$ найден коэффициент отражения от конца волноводного тракта $\Gamma_{\text{к}}$ и оценён КПД системы по рассогласованию.

По измеренному периоду определена длина волны в волноводе ${\lambda^{\text{п}}_{\text{в}}} = 4{,}8$~см, которой соответствует длина волны в свободном пространстве ${\lambda^{\text{п}}_{\text{с}}} \approx 3{,}32$~см. КНД определён двумя независимыми способами~--- по зависимости $I(\Delta X)$ и по максимуму зависимости $\widetilde{\Gamma}(\Delta X)$. Получены $D_2 \approx 43{,}7$ ($16{,}4$~дБ) и $D_3 \approx 36{,}8$ ($15{,}7$~дБ). Они отличаются примерно на $17\%$, что объясняется дискретностью шага перемещения, неточностью определения экстремумов и приближением, в котором пренебрегается квадратичными по $\Gamma$ членами. Проверка показала, что эти члены действительно малы, поэтому использованное приближение корректно. Оценка свободной длины волны по периодичности осцилляций во втором задании ($34$~мм) согласуется с результатом первого задания. Средний экспериментальный КНД ниже теоретической оценки $D_{\text{теор}} \approx 142{,}3$ примерно на $71{,}7\%$, что связано с неоднородностью амплитуды и фазы поля на раскрыве и неидеальностью установки.
